\documentclass[12pt]{article}
\usepackage[T1]{fontenc}
\usepackage[utf8]{inputenc}
\usepackage{lmodern}
\usepackage{textcomp}
\usepackage{enumitem}
\usepackage{geometry}
\geometry{margin=1in}
\begin{document}
\begin{center}
Answer Key - Health Sciences - Undergraduate Level Assessment 10 \
\small (Answers are ordered exactly as in the question paper.)
\end{center}

\section*{Multiple Choice}
\begin{enumerate}[label=\arabic*.]
\item \textbf{Answer:} D
\\ \textit{Options:} A) bilateral contraction of antigravity limb muscles.; B) bilateral contraction of limb musculature.; C) contraction of ipsilateral limb musculature.; D) contraction of contralateral limb musculature.
\\ \textit{Note:} The primary motor cortex is responsible for initiating voluntary movements and directly influences the contraction of muscles on the opposite side of the body (contralateral) due to the crossing of motor pathways in the brainstem.
\item \textbf{Answer:} D
\\ \textit{Options:} A) Moving to be close to their children; B) Moving to Arizona or Florida if they could; C) Residing in a very nice nursing home; D) Staying in their own home
\\ \textit{Note:} Most older adults prefer staying in their own home because it provides them with a sense of familiarity, independence, and comfort, which are crucial for their emotional well-being and quality of life.
\item \textbf{Answer:} A
\\ \textit{Options:} A) Amplify by persistent infection of midgets, mosquitoes, sand flies, ticks; B) Predominantly mosquito borne; C) Genetic recombination or reassortment; D) Spread to humans via bites of mice and rats
\\ \textit{Note:} The bunyavirus family is characterized by its ability to replicate and persist within arthropod vectors such as mosquitoes and ticks, which play a crucial role in their transmission and life cycle.
\item \textbf{Answer:} B
\\ \textit{Options:} A) Single nucleotide polymorphisms (SNPs); B) Microsatellites; C) Minisatellites; D) Satellites
\\ \textit{Note:} Microsatellites are defined as short, repeating sequences of 2 to 4 base pairs of DNA, which makes them the correct answer to the question about repeat core sequences.
\item \textbf{Answer:} A
\\ \textit{Options:} A) prophase I; B) metaphase I; C) prophase II; D) metaphase II
\\ \textit{Note:} Prophase I is the stage of meiosis where homologous chromosomes pair up and engage in crossing over, allowing for genetic recombination and increased genetic diversity.
\end{enumerate}

\section*{True / False}
\begin{enumerate}[label=\arabic*.]
\setcounter{enumi}{5}
\item \textbf{Answer:} False
\\ \textit{Options:} A) True; B) False
\\ \textit{Note:} Monosomy is not the most frequently observed chromosome abnormality in first trimester spontaneous miscarriages; rather, trisomy, particularly trisomy 16, is more commonly identified in these cases.
\item \textbf{Answer:} True
\\ \textit{Options:} A) True; B) False
\\ \textit{Note:} Research indicates that many girls tend to explore their own bodies and develop an understanding of masturbation through personal experience rather than relying on external sources such as education or media.
\item \textbf{Answer:} False
\\ \textit{Options:} A) True; B) False
\\ \textit{Note:} Cytotoxic T cells are primarily activated by recognizing and binding to antigen-presenting cells that display viral peptides on their surface, rather than through interaction with antibodies in the bloodstream.
\item \textbf{Answer:} True
\\ \textit{Options:} A) True; B) False
\\ \textit{Note:} Herpes viruses primarily cause conditions such as cold sores and genital herpes, and they are not associated with causing infantile paralysis, which is typically linked to viruses like poliovirus.
\item \textbf{Answer:} True
\\ \textit{Options:} A) True; B) False
\\ \textit{Note:} Sterilization is the most common method of birth control for married couples in the U.S. because it provides a permanent solution to prevent pregnancies, making it a preferred choice for those who have completed their families.
\end{enumerate}

\section*{Long Form Response}
\begin{enumerate}[label=\arabic*.]
\setcounter{enumi}{10}
\item \textbf{Answer:} Proprioceptive information plays a crucial role in the function of the temporomandibular joint (TMJ) by providing the central nervous system with feedback about the position and movement of the jaw. This sensory information is primarily conveyed through the masseteric and auriculotemporal nerves, which innervate the masseter muscle and the TMJ, respectively. The masseteric nerve transmits proprioceptive signals related to the muscle's contraction and tension, while the auriculotemporal nerve carries sensory information from the joint itself. Together, these nerves facilitate coordinated jaw movements and help maintain proper occlusion and joint stability, which are essential for effective mastication and overall oral health. Understanding the proprioceptive pathways involved in TMJ function underscores the importance of these nerves in the regulation of masticatory dynamics and the prevention of dysfunction.
\\ \textit{Note:} correct because it accurately describes how proprioceptive information from the masseteric and auriculotemporal nerves provides essential sensory feedback to the central nervous system regarding the position and movement of the jaw, which is vital for the proper function of the temporomandibular joint.
\item \textbf{Answer:} Ovulation is a critical phase of the typical 28-day menstrual cycle, occurring around day 14. This process is initiated by a surge in luteinizing hormone (LH) triggered by elevated levels of estrogen produced by developing follicles. Physiologically, the leading follicle releases an egg, which is then available for fertilization. Factors such as hormonal imbalances, stress, and lifestyle choices can influence the timing and regularity of ovulation. Understanding these elements is essential in health sciences, as they play a significant role in reproductive health and fertility management.
\\ \textit{Note:} The answer correctly identifies ovulation as occurring around day 14 of the menstrual cycle, explains the hormonal regulation involving a surge in luteinizing hormone and estrogen, describes the physiological release of the egg, and acknowledges external factors that can impact ovulation, thus providing a comprehensive overview of the process within the specified context.
\end{enumerate}

\vfill
\noindent\textit{End of Answer Key}
\end{document}